% © 2026 Ing. David Jaroš — CC BY-NC-ND 4.0
%
% This work is licensed under a Creative Commons Attribution-NonCommercial-NoDerivatives
% 4.0 International License (CC BY-NC-ND 4.0).
%
% License History: Earlier drafts (up to v0.3) were released under CC BY 4.0.
% From v0.4 onward, all material is released under CC BY-NC-ND 4.0 to protect
% the integrity of the theoretical work during ongoing academic development.
%
% See LICENSE.md for full license text.

\ifdefined\INCLUDEMODE
\else
  \documentclass[12pt,a4paper]{article}
  \usepackage[a4paper, margin=2.5cm]{geometry}
  \usepackage{amsmath, amssymb, amsthm}
  \usepackage{mathtools}
  \usepackage{hyperref}
  \usepackage{xcolor}
  \usepackage{mdframed}
  \usepackage{booktabs}

  \newtheorem{theorem}{Theorem}[section]
  \newtheorem{lemma}[theorem]{Lemma}
  \newtheorem{proposition}[theorem]{Proposition}
  \newtheorem{conjecture}[theorem]{Conjecture}
  \theoremstyle{definition}
  \newtheorem{definition}[theorem]{Definition}
  \theoremstyle{remark}
  \newtheorem{remark}[theorem]{Remark}

  \newmdenv[backgroundcolor=yellow!20, linecolor=orange, linewidth=1.5pt]{gapbox}
  \newmdenv[backgroundcolor=green!15, linecolor=green!60!black, linewidth=1.5pt]{resultbox}
  \newmdenv[backgroundcolor=red!10, linecolor=red!60, linewidth=1.5pt]{deadendbox}
  \newmdenv[backgroundcolor=blue!8, linecolor=blue!50, linewidth=1.5pt]{insightbox}

  \title{\textbf{Electron Mass --- Step 3: Modular Period Approach}\\
         \large Can the Period of a Modular Form Fix $m_e$?}
  \author{David Jaro\v{s}}
  \date{2026-03-06}

  \begin{document}
  \maketitle

  \begin{abstract}
  We explore whether the electron mass $m_e$ is related to the period
  $\Omega(f)$ of the Hecke newform $f_1$ encoding the electron generation.
  Periods of modular forms are transcendental but computable numbers that
  encode the arithmetic of the modular curve.  We test: is
  $m_e = C\,\Omega(f_1)$ for some natural constant $C$ with a UBT
  interpretation?  We find that this approach requires a dimensional
  bridge between the period (a length in modular space) and $m_e$
  (a mass), which is not present in the current UBT framework.
  Gap~Y2 remains open.
  \end{abstract}
\fi

%──────────────────────────────────────────────────────────────────────────────
\section{Modular Periods as Transcendental Constants}
\label{sec:periods}
%──────────────────────────────────────────────────────────────────────────────

For a weight-2 newform $f \in S_2(\Gamma_0(N))$ with Fourier expansion
$f(\tau) = \sum_{n\geq 1} a_n\,q^n$ ($q = e^{2\pi i\tau}$), the
\emph{period lattice} is spanned by:
\begin{equation}
  \Omega^+(f) \;=\; \int_0^{i\infty} f(\tau)\,d\tau, \qquad
  \Omega^-(f) \;=\; \int_0^{1/2} f(\tau)\,d\tau.
\end{equation}
The real period $\Omega^+(f) = -i\,L(1,f)$ is computable from the
$L$-function.

\begin{insightbox}
\textbf{Connection to the Hecke conjecture:}
If $f_1$ is the electron-generation form (Step~1) and the Hecke
conjecture holds, then $\Omega^+(f_1)$ is a natural modular-arithmetic
quantity associated with $m_e$.
\end{insightbox}

%──────────────────────────────────────────────────────────────────────────────
\section{Dimensional Analysis}
\label{sec:dimensional}
%──────────────────────────────────────────────────────────────────────────────

The period $\Omega^+(f_1)$ is a dimensionless complex number (in the
convention $d\tau = i\,dt$, it has dimensions of area in $\tau$-space,
but in the standard normalisation it is a pure number).  The electron
mass $m_e$ has dimensions of energy.

To relate them, one needs a dimensional bridge:
\begin{equation}
  m_e c^2 \;=\; \frac{\hbar\,c}{R_\psi} \times F(\Omega^+(f_1)),
\end{equation}
where $F$ is some function and $R_\psi$ provides the dimensional factor.
But $R_\psi$ is calibrated to $m_e$, making this circular.

An alternative: use the Planck scale as the bridge:
\begin{equation}
  m_e \;=\; m_P \times F(\Omega^+(f_1)),
\end{equation}
where $m_P = \sqrt{\hbar c/G} \approx 2.18 \times 10^{-8}\;\mathrm{kg}$.
Then $m_e/m_P \approx 4.19 \times 10^{-23}$.  This requires
$F(\Omega^+(f_1)) \approx 4 \times 10^{-23}$, an extremely small number
that cannot be a simple modular period value.

%──────────────────────────────────────────────────────────────────────────────
\section{Numerical Test}
\label{sec:numerical}
%──────────────────────────────────────────────────────────────────────────────

For the LMFDB label \texttt{38.2.a.a} (pending David's SageMath run),
the numerical value $L(1,f_1) \approx O(1)$ gives:
\begin{equation}
  |\Omega^+(f_1)| = |L(1,f_1)| \;\sim\; O(1).
\end{equation}
To give $m_e = 0.511\;\mathrm{MeV}$, the constant $C$ must carry the
mass dimension:
\begin{equation}
  m_e \;=\; C \;=\; \frac{\hbar}{c\,R_\psi} \times O(1),
\end{equation}
which is again calibrated to $R_\psi \propto 1/m_e$.

%──────────────────────────────────────────────────────────────────────────────
\section{A Speculative Arithmetic Connection}
\label{sec:speculative}
%──────────────────────────────────────────────────────────────────────────────

There is a deep conjecture (Birch and Swinnerton-Dyer, for $k=2$) that
$L(1,f)$ is related to the N\'eron-Tate heights and ranks of associated
elliptic curves.  For the electron-generation form $f_1$ (associated
to an elliptic curve $E$ of conductor 38):
\begin{equation}
  L(1,f_1) \;=\; \frac{\Omega(E) \times \prod c_p \times |\Sha(E)|}
                       {|E(\mathbb{Q})_{\mathrm{tors}}|^2},
\end{equation}
where $\Omega(E)$ is the real period, $c_p$ are Tamagawa numbers,
$\Sha(E)$ is the Tate-Shafarevich group, and $E(\mathbb{Q})_{\mathrm{tors}}$
is the torsion subgroup.  All of these are computable but require the
explicit elliptic curve.

\begin{conjecture}[Speculative: $m_e$ from BSD]
\label{conj:BSD}
The electron mass is given by:
\begin{equation}
  m_e c^2 \;=\; \frac{\hbar c}{\ell_P} \times \frac{\Omega(E_e)}{\text{vol}(X_0(38))},
\end{equation}
where $E_e$ is the elliptic curve of conductor 38 associated with $f_1$,
$\ell_P$ is the Planck length, and vol$(X_0(38))$ normalises to the
modular curve area.
\end{conjecture}

\begin{remark}
Conjecture~\ref{conj:BSD} is highly speculative and not yet testable
without the explicit identification of $E_e$ (pending LMFDB).  It is
presented as a research direction, not a result.
\end{remark}

%──────────────────────────────────────────────────────────────────────────────
\section{Conclusion}
\label{sec:conclusion}
%──────────────────────────────────────────────────────────────────────────────

\begin{gapbox}
\textbf{Status: OPEN.  Gap~Y2 not closed.}

The modular period approach requires a dimensional bridge between
the dimensionless period $\Omega^+(f_1) \sim O(1)$ and the electron mass.
No natural bridge exists in current UBT without either:
(a) Calibrating $R_\psi$ to $m_e$ (circular), or
(b) Using the Planck scale (gives $m_e/m_P \sim 10^{-23}$, which modular
    periods cannot naturally produce).

\textbf{Honest statement:} The three approaches to fixing $m_e$ in Steps~1--3
are:
\begin{itemize}
  \item Step~1: Conditional on Hecke --- reduces from 3 free to 1 free
        parameter.  This is real progress.
  \item Step~2: Geometric $R_\psi$ --- dead end.
  \item Step~3: Modular period --- open/speculative.
\end{itemize}

Gap~Y1 (derive $y_e$ from $S[\Theta]$) and Gap~Y2 (derive $v$ without
$m_e$ input) remain open.  The electron mass scale is the fundamental
free parameter of the UBT lepton sector (conditional on Hecke).
\end{gapbox}

\ifdefined\INCLUDEMODE
\else
  \end{document}
\fi
